% =====================================================================
%  Betriebsanweisung Tauchsäge Festool TS 55 REBQ mit Führungsschiene
%  Eigenes PDF: Betriebsanweisung_Kreissaege.tex, Seite in der Einweisung
%  über \BAseite. Layout: fablab-document/fablab-ba.sty
%
%  Lizenz-Hinweis: Text vollständig selbst formuliert, KEINE Passagen und
%  KEINE Abbildungen aus der Festool-Betriebsanleitung übernehmen! Sicherheits-
%  regeln als solche (Fakten) sind frei, der Wortlaut und die Bilder von
%  Festool sind urheberrechtlich geschützt. Nur ISO-7010-Symbole verwenden.
% =====================================================================
\BAsetup{
  typ=maschine,
  titel={Tauchsäge Festool TS 55 REBQ mit Führungsschiene},
  kurztitel={Tauchsäge},
  taetigkeit={Sägen mit Führungsschiene, Sägeblattwechsel, Reinigen},
  nummer=BA-HK-01,
  stand=29.09.2026,                   % Datum der Freigabe
  verantwortlich={Name der verantwortlichen Person},
  signalwort=WARNUNG,
  schrift=\footnotesize,
}

\BAkopf

\BAblock{ANWENDUNGSBEREICH}{M002}{%
  \begin{itemize}
    \item Sägen von Plattenmaterial, Holz, Kunststoff und (mit Aluminium-Sägeblatt) Aluminium.
    \item Nur mit Führungsschiene, freihändiges Sägen ist nicht erlaubt. Nicht während der normalen Öffnungszeiten.
    \item Benutzung nur nach unterschriebener Einweisung; Betriebsanleitung des Herstellers beachten.
  \end{itemize}}

\BAblock{GEFAHREN FÜR MENSCH UND UMWELT}{W022,W024,W012}{%
  \begin{itemize}
    \item Schnittverletzungen, vor allem unter dem Werkstück (dort schützt die Haube nicht).
    \item \textbf{Rückschlag:} Verkantet das Sägeblatt, kann die Säge unkontrolliert zur Bedienperson springen.
    \item Wegfliegende Teile und Späne; Stromschlag bei angesägtem Netzkabel; Lärm.
    \item Gesundheitsschäden durch Staub; Hartholzstaub kann Krebs erzeugen.
  \end{itemize}}

\BAblock{SCHUTZMASSNAHMEN UND VERHALTENSREGELN}{M003,M004,M017,P028}{%
  \begin{itemize}
    \item Vor jeder Benutzung Säge, Kabel, Blatt und Schutzhaube prüfen; defekte Säge nicht benutzen.
    \item Gehörschutz tragen; Schutzbrille bei Metall; Staubmaske bei viel Staub (Hartholz). Anliegende Kleidung, Haare zusammenbinden. \textbf{Beim Sägen keine Handschuhe} (Einzugsgefahr).
    \item Werkstück und Führungsschiene fest spannen. Immer Absaugung (Festool-Sauger, Staubklasse M) anschließen; bei Hartholzstaub Arbeit abbrechen bzw.\ ins Freie verlegen.
    \item Säge mit \textbf{beiden Händen} führen, seitlich der Sägeebene stehen. Nie unter das Werkstück greifen.
    \item Schnitttiefe anpassen: Blatt ragt nur wenig unter dem Werkstück heraus, im Arbeitstisch höchstens 1\,mm tief sägen.
    \item Säge nur eingeschaltet ins Material führen, nur vorwärts schieben. Bei klemmendem Blatt Schalter loslassen und Säge ruhig halten, bis das Blatt steht.
    \item Kabel aus dem Schnittbereich fernhalten; bei Kabelschaden sofort Stecker ziehen.
    \item Aluminium: nur Aluminium-Sägeblatt, Spanflugschutz schließen. Nur passende Original-Sägeblätter verwenden.
  \end{itemize}}

\BAblock{VERHALTEN BEI STÖRUNGEN UND IM GEFAHRENFALL}{W001,M006}{%
  \begin{itemize}
    \item Säge ausschalten, Stecker ziehen, Betreuer informieren; Schaden und Hergang per Mail an \texttt{fablab-aktive@fablab.fau.de} melden.
    \item Entstehungsbrand: Säge vom Netz trennen, löschen. \textbf{Feuerwehr: 112}
  \end{itemize}}

\BAblock{VERHALTEN BEI UNFÄLLEN – ERSTE HILFE}{E003}{%
  \begin{itemize}
    \item Säge ausschalten, Stecker ziehen. Betreuer/Ersthelfer verständigen, Wunden versorgen.
    \item Jeden Unfall melden und im Verbandbuch eintragen.
  \end{itemize}\vspace{1.5mm}
  \textbf{Notruf: 112}\qquad FabLab: 09131\,/\,85\,28013}

\BAletzterblock{INSTANDHALTUNG, ENTSORGUNG}{M021}{%
  \begin{itemize}
    \item Sägeblattwechsel und Wartung nur durch eingewiesene Personen, mit gezogenem Stecker und Schutzhandschuhen.
    \item Säge und Motorlüftung regelmäßig von Staub reinigen. Elektrische Prüfung nach DGUV Vorschrift 3 gemäß Prüffrist.
  \end{itemize}}
